% RESEARCH_DRAFT_SHA256: 5f026ed23665316c25df86cc5f56378b6e2f563d9cc1fad19e8dac5e8e062ac7
% Layout and source-reference wrappers expanded for the web renderer.
\documentclass[11pt]{article}
\usepackage[margin=25mm]{geometry}
\usepackage{amsthm}
\usepackage{fontspec,unicode-math}
\setmainfont{Latin Modern Roman}
\setmathfont{Latin Modern Math}
\usepackage{amsmath,amssymb,mathtools,xurl,hyperref}
\hypersetup{colorlinks=true,urlcolor=blue,linkcolor=blue}
\setcounter{secnumdepth}{0}
\setlength{\parindent}{0pt}
\setlength{\parskip}{5pt}
\setlength{\emergencystretch}{3em}
\newcommand{\sref}[2]{\hyperlink{source:#1:#2}{[S#2]}}
\newcommand{\cataloguescope}[1]{\paragraph{Source record.}#1}
\newcommand{\floor}[1]{\left\lfloor#1\right\rfloor}
\newtheorem{proposition}{Proposition}
\newtheorem{lemma}[proposition]{Lemma}
\begin{document}
\paragraph{Research manuscript.} OpenAI \texttt{gpt-6-astra}, reasoning \texttt{ultra}; 11 October 2026.
\paragraph{Preserved source body.} The complete catalogue mathematical text and sources follow with layout and source links expanded, including source qualifications and its historical preparation notes. New research begins at the section headed Research attempt.

\section{Efficient Optimization of Neural-Network Weights}\label{statement}

\subsection{Definitions and mathematical statement}

\paragraph{Shared notation.}Standard number systems are $\mathbb N,\mathbb Z,\mathbb Q,\mathbb R,\mathbb C$. Specialized notation and index ranges follow the cited source and the definitions below.


Fixing a neural-network architecture fixes its graph and activations, while weights are the trainable numerical parameters. Global optimization minimizes the specified training objective over all admissible weights.
\paragraph{Record scope and qualifications.}
The catalogue title is retained as metadata. Named certificates, contracts, bounds or frameworks in that title are not asserted as results of the original source extract.
The loss, data encoding, exact versus approximate optimum, weight constraints, and meaning of reasonable time are not specified; fixed architecture does not automatically mean fixed input size.

Source-record limitation: The loss, input representation, admissible weights, approximation criterion, and reasonable-time requirement are unspecified.
\paragraph{Statement recorded in the supplied source.}
Given a fixed neural-network architecture, can globally optimal weights be found in reasonable computational time? \href{https://doi.org/10.4171/owr/2025/22}{[S2]}
\paragraph{Source status.}
The supplied catalogue labels this source record open; that metadata was checked against the supplied snapshot on 2026-10-11. The label is preserved for provenance and is not an independent verification of present open status or mathematical correctness. The catalogue literature assessment is dated 2026-08-22.
\subsection{Short English statement}
Determine when globally best weights of a fixed neural network can be found efficiently.
\subsection{Sources}
\begin{itemize}
\item[S1] Ulam AI and the UnsolvedMath Contributors, UnsolvedMath: A Curated Collection of Open Mathematics Problems, catalogue entry. Catalogue curation and metadata: CC BY 4.0. The source extract is retained; definitions, English summaries and qualifications are editorial additions.
\url{https://huggingface.co/datasets/ulamai/UnsolvedMath}.
\item[S2] Computational Multiscale Methods (2025), Oberwolfach Reports, DOI 10.4171/owr/2025/22, pp. 1079--1162
\url{https://doi.org/10.4171/owr/2025/22}.
\item[S3] Abrahamsen--Bouritsas--Dughmi--Miltzow, training complexity
\url{https://papers.nips.cc/paper_files/paper/2023/hash/71c31ebf577ffdad5f4a74156daad518-Abstract-Conference.html}.
\item Source-specific Oberwolfach Reports licences, including CC BY-SA 4.0 where applicable, are retained; see THIRD\_PARTY\_NOTICES.md.
\end{itemize}
\paragraph{Preparation and review.}
The original source and UnsolvedMath contributors retain attribution for the source extract. Codex (GPT-6; exact serving revision not exposed) assisted with definitions, English summaries, scope notes and formatting. No new human mathematical authorship or independently verified solution is asserted. Catalogue identity checked against the complete supplied registry and current repository. Distinct upstream records, including repeated questions and broad or damaged extracts, are retained for complete coverage of the source CS category. A repeated record is not a new independent conjecture. Complete mathematical-equivalence and current-literature review has not been performed. Preparation and final review follow contribution prompts C and D, with source gaps disclosed. Structural and rendering checks do not establish a new theorem or open status.
\label{end:statement}

\section{Research attempt}
\subsection{Literature, source context and scope}
The full original catalogue body and source qualifications are preserved above. The relevant primary workshop passages were read in the cached publisher report. Historical claims and subsequent attributed results are distinguished from the mathematical work below. A restricted example, standard lemma or known theorem is not counted as new progress on the entire original target. Missing assumptions and repeated or rhetorical extracts are not silently turned into different solved problems.

The report asks about efficient weight computation even when an approximation architecture is known, and describes a hierarchical approach under a stability hypothesis. It does not specify one loss or a uniform complexity model for all architectures. This attempt studies exact squared-loss training of one scalar ReLU and then locates the failure in extending gradient-based reasoning even to that model.
\subsection{A complete global algorithm for a single scalar weight}
For rational data $(x_i,y_i)_{i=1}^N$, $N\ge1$, consider the fixed architecture $f_w(x)=\max(0,wx)$ with one real weight and no bias, and loss $L(w)=\sum_i(f_w(x_i)-y_i)^2$. Set
\[
 z_i^+=\max(x_i,0),\quad z_i^-=\min(x_i,0),\quad
 A_\pm=\sum_i(z_i^\pm)^2,\quad B_\pm=\sum_i y_i z_i^\pm,\quad C=\sum_i y_i^2.
\]
\begin{proposition}
A global optimum is found by evaluating just two candidates:
\[
 w_+=\begin{cases}\max(0,B_+/A_+),&A_+>0,\\0,&A_+=0,\end{cases}
 \quad
 w_-=\begin{cases}\min(0,B_-/A_-),&A_->0,\\0,&A_-=0.\end{cases}
\]
Choose the one of smaller loss. This uses $O(N)$ rational arithmetic operations and polynomial bit complexity in the rational input length.
\end{proposition}
\begin{proof}
On $w\ge0$, the prediction is $wz_i^+$, and on $w\le0$ it is $wz_i^-$. Thus the loss on each closed half-line is $A_\pm w^2-2B_\pm w+C$. If $A_\pm>0$, completing the square gives its unconstrained minimizer $B_\pm/A_\pm$, clipped to the appropriate half-line as displayed. If $A_\pm=0$, every $z_i^\pm=0$, hence also $B_\pm=0$ and the loss is constantly $C$ there. The two half-lines cover every real weight, so comparison gives a global optimum. Computing the sums is linear in $N$. With rational numerators and denominators in binary, naive common-denominator additions and multiplications give intermediate bit lengths bounded by a polynomial in the total input length, and ordinary integer arithmetic then gives polynomial bit cost.
\end{proof}
\subsection{A failed stationary-point route}
For the single data point $(x,y)=(1,1)$, every $w<0$ has loss one and derivative zero, while $w=1$ has loss zero. Exact gradient descent started at any negative weight remains there, independently of its step sizes, since each computed gradient is zero. This is a suboptimal open plateau, not a nondifferentiability issue at zero. Thus a certificate of stationarity is insufficient even where the global optimum is cheaply computable.
\section{Verification and remaining gaps}
Exact rational checks compare the two-candidate solution against each half-line quadratic for several datasets, and verify the dead-ReLU plateau. For more weights the activation regions depend jointly on parameters; enumerating them may become expensive and multilayer loss is no longer quadratic on each region. No general efficient optimal-weight algorithm or complexity lower bound for the source's broad problem follows from this one-weight calculation.

\section{References and source attribution}
Primary workshop report: \url{https://doi.org/10.4171/owr/2025/22}. Relevant passages read: OWR 2025/22, printed pages1090--1091, Feischl hierarchical training.

All retrievals were checked on 11 October 2026. Original question attribution remains in the preserved Sources section. External theorems are attributed and not claimed as new discoveries or fully verified merely by citation.


\paragraph{Integration review.} Codex (GPT-6; exact serving revision not exposed) reviewed the written arguments and expanded layout and source-reference wrappers. Repository metadata and the production web renderer were checked using Python 3.13 and JavaScript with MathJax 3.2.2. This model-assisted review is not independent human proof verification; no PDF compilation or formal proof checking is asserted.

\section{Completion Estimate}
\noindent\textbf{COMPLETION: 0\%}
No substantiated new progress toward the entire original target is claimed. The worked results and failed route are delimited above; a known result, illustrative model, or auxiliary obstruction is not counted as original-target completion.

\section{Final status and attribution}
\textbf{Attempt status: unresolved.} This describes this research attempt, not a new ruling on the source catalogue's status or on external literature claims.
The mathematics and this self-audit were generated by OpenAI \texttt{gpt-6-astra}, reasoning \texttt{ultra}, on 11 October 2026. Codex (GPT-6; exact serving revision not exposed) orchestrated the batch, read and reviewed the written mathematical arguments, independently reran the finite checks, and contributed formatting/repository review. No human mathematical authorship, independent human proof checking, or formal verification is claimed. The historical preparation paragraphs in the reproduced catalogue body concern that body only.

Contribution prompt A was applied to the complete definition and supplied source record, followed by the adversarial checks of prompt B. Local reproduction files are \texttt{batch026/checks.py}, \texttt{computations.json}, and the per-ID audit JSON. They were not public links at preparation time. Numerical checks are finite and do not prove unrestricted results. No LaTeX engine was found locally; compilation was not performed by this research agent.

\end{document}
